\section{Literature Review}

\subsection{Traditional and Data-Mining Approaches to Tax Fraud Detection}

According to Castellón, tax authorities have traditionally tackled tax fraud through two approaches: auditors experience and rule-based systems~\cite{gonzalez2013characterization}. The former approach consists of randomly selecting tax declarations and auditing them based on the experience, intuition, and domain knowledge of the tax auditors. The latter uses methodologies based on rule-based systems. A rule-based system, as described by Baesens et al., is often implemented in the form of a set of if-then rules that detect fraud cases~\cite{baesens2015fraud}. These rules are developed by a burdensome process, where auditors, after a through review, detect a tax fraud case, generalize its characteristics and include a rule into the tax fraud knowledge base. Once applied to incoming cases it triggers an alert or signal when fraud is detected. However, these traditional techniques have two important disadvantages: (1) They rely mainly on past experiences, therefore are unable to discover new fraud mechanisms by themselves. (2) The subjective judgment of experts makes the knowledge bases for rule-based systems expensive to build, maintain and update~\cite{baesens2015fraud}. A more recent approach to detect tax fraud is through the use of data mining techniques, which provide mechanisms to extract and generate knowledge from substantial volumes of data to support the detection of fraudulent behavior, improving the use of resources~\cite{gonzalez2013characterization}. Also, this is one of the best established applications of data mining in both government and industry. Tax authorities in most countries have frequently adopted data mining techniques to assist them in identifying taxpayers who evade obligations~\cite{gonzalez2013characterization}. There are works, in the literature, using supervised learning techniques because of the use of labeled data or audit-assisted data.

\subsection{Machine Learning for Tax Risk and Financial Anomaly Detection}

The introduction of machine learning technology has transformed tasks such as anomaly detection in corporate financial data, tax risk prediction, and compliance auditing from static analysis to dynamic modeling~\cite{olorunyomi2022dynamic,bello2024implementing}. Research based on different models and algorithmic frameworks has gradually expanded, building efficient risk early warning systems and methods for detecting tax-related behavior, which not only provide strong support for financial supervision but also offer intelligent pathways for enterprises to optimize tax compliance strategies~\cite{yang2025predictive}.

Financial anomaly detection is an important application of machine learning technology. Researchers have built efficient anomaly identification models to analyze and warn of potential abnormal behavior in corporate financial data. Wang et al. developed an XGBoost ensemble model to identify abnormal tax burden levels among enterprises, significantly enhancing classification performance under imbalanced sample conditions~\cite{wang2022hybrid}. Mojahedi et al. compared the performance of Logistic Regression, Decision Tree, and Support Vector Machine (SVM) on tax-related datasets, finding that Logistic Regression performed stably when dealing with low-dimensional data, while SVM was more suitable for anomaly detection in nonlinear scenarios~\cite{mojahedi2022towards}. Zhang et al. designed a triple feature selection mechanism based on the Random Forest model, which effectively improved detection accuracy and stability under multi-dimensional data conditions~\cite{zhang2023three}. Liu et al. further proposed a distribution-transformation-based Random Forest approach, demonstrating strong performance in detecting abnormal behaviors in enterprises' inter-period financial fluctuations~\cite{liu2025self}. Meanwhile, George et al., through empirical research, found that traditional machine learning models still hold significant application value in identifying financial fraud risks, particularly in enterprise scenarios with moderately sized data samples and clearly defined feature variables~\cite{george2024predictive}.

\subsection{Unsupervised Approaches}

While most studies addressing tax fraud detection employ supervised learning, the lack of historically labeled data limits such studies in most situations since manual auditing and labeling data is expensive and time-consuming. As a result, they proposed using unsupervised learning~\cite{roux2018tax}. The research of de Roux et al. was divided into three phases, which are clustering similarly valued tax declarations; adjusting probability distribution to each cluster tax base, and detecting suspicious activities by checking the quantiles of cluster adjusted distribution. Despite the fact that their model produced promising results, there was the limitation of the limited amount of data, and a low number of variables needed to deeply describe cases. The most worrying issue was that the model's accuracy would not be evaluated because the data was not labeled, which raises questions about its trustworthiness~\cite{murorunkwere2022fraud}.

\subsection{Anomaly Detection Methods}

Chandola et al.~\cite{chandola2009anomaly} provide the standard taxonomy; the profiles sought here are point anomalies in a multivariate financial space. Anomaly is defined compared to the chosen model normality~\cite{aggarwal2017outlier}. Isolation Forest~\cite{liu2008isolation,liu2012isolation} isolates records by splitting them using randomly selected feature. Those which require fewer splits are considered as anomalous. The Local Outlier Factor~\cite{breunig2000lof} instead compares a record's local density with that of its nearest neighbours. Comparative studies find no universally superior detector: Goldstein and Uchida~\cite{goldstein2016comparative} found that relative performance depends on whether anomalies are global or local, Campos et al.~\cite{campos2016evaluation} documented strong sensitivity to parameter choice, and Zimek et al.~\cite{zimek2012survey} analyse neighbourhood-based detectors in high-dimensional spaces where distance contrast deteriorates.

\subsection{Explainable Artificial Intelligence}

The use of Explainable Artificial Intelligence (XAI) methods, including SHAP (SHapley Additive Explanations) and LIME (Local Interpretable Model-Agnostic Explanations), is becoming more widespread to explain the predictions of the model, so that they can be regulated and trusted by the auditors~\cite{chen2024explainable}.

\subsection{Research Gap and Synthesis}

The literature shows that most tax fraud detection studies rely on supervised learning, which requires labelled audit outcomes that are rarely available, while the unsupervised studies that do exist score every taxpayer against the whole population and offer little evaluation or explanation of their results. Comparative studies also show that no anomaly detector is universally superior and that the reference population and feature space strongly influence what is flagged. This research addresses these gaps by segmenting taxpayers into peer groups before applying Isolation Forest, evaluating the results against injected anomalies and a supervised comparator, and explaining both segment membership and anomaly flags with SHAP, as set out in the research questions in Chapter 1.
